% ==============================================================================
% 11_evaluation.tex — UE SIS589
% ==============================================================================
\chapter{Évaluation et Grille de Compétences}
\label{chap:evaluation}

\epigraph{\itshape « L'examen n'est pas la fin de l'apprentissage~;
  c'en est la mesure. »}
         {--- Principe pédagogique ENSP}

% ==============================================================================
\section{Objectifs et structure de l'évaluation}
\label{sec:eval-structure}
% ==============================================================================

L'évaluation de l'UE SIS589 porte sur trois dimensions~:

\begin{enumerate}
  \item \textbf{Connaissances théoriques~:} maîtrise des concepts,
        définitions, normes et référentiels.
  \item \textbf{Compétences techniques~:} capacité à analyser des
        artefacts, rédiger des règles, interpréter des journaux,
        conduire une investigation.
  \item \textbf{Compétences professionnelles~:} raisonnement en
        situation, documentation, communication.
\end{enumerate}

\textbf{Répartition des notes~:}
\begin{itemize}
  \item Contrôle continu (QCM en ligne après chaque partie)~: 20\,\%
  \item Rapports de laboratoires (Lab~1, 2, 3)~: 40\,\%
  \item Examen final (4~heures)~: 40\,\%
\end{itemize}

% ==============================================================================
\section{Grille de compétences}
\label{sec:competences}
% ==============================================================================

\begin{table}[H]
\centering\small
\caption{Grille de compétences SIS589 --- niveaux de maîtrise}
\label{tab:competences}
\rowcolors{2}{TableauPair}{TableauImpair}
\begin{tabular}{p{5cm}p{2.5cm}p{6cm}}
\toprule
\tableauHeader{Compétence} & \tableauHeader{Niveau attendu} &
\tableauHeader{Indicateurs de maîtrise} \\
\midrule
Architecturer un SOC minimal &
Opérationnel &
Choisir les composants, les positionner, justifier le choix \\
Configurer un SIEM (Wazuh + ELK) &
Pratique &
Déployer agents, pipeline Logstash, tableau de bord \\
Rédiger une règle de détection &
Pratique &
Syntaxe Wazuh ou Sigma correcte, mapping MITRE \\
Analyser des journaux multi-sources &
Maîtrise &
Identifier les événements pertinents, construire la timeline \\
Conduire une acquisition forensique &
Pratique &
dd / FTK Imager, hash, PV de collecte conforme ISO~27037 \\
Analyser un dump mémoire &
Opérationnel &
Commandes Volatility~3 essentielles, interprétation \\
Analyser un PCAP &
Opérationnel &
Filtres Wireshark, identification des flux suspects \\
Piloter une réponse à incident &
Opérationnel &
Cycle ISO~27035, qualification gravité, confinement \\
Rédiger un rapport d'incident &
Maîtrise &
Structure en 8~sections, niveau exécutif + technique \\
Utiliser MISP / CTI &
Notionnel &
Rechercher un IoC, créer un événement, exporter pour SIEM \\
Appliquer le cadre légal camerounais &
Notionnel &
Identifier les infractions, obligations ANTIC, chaîne de custody \\
Conduire un RETEX &
Opérationnel &
6~sections, 5~Pourquoi, plan de remédiation priorisé \\
\bottomrule
\end{tabular}
\end{table}

\textbf{Niveaux~:}
\textit{Notionnel}~= connaît le concept et peut l'expliquer~;
\textit{Opérationnel}~= peut appliquer en situation guidée~;
\textit{Pratique}~= applique de manière autonome~;
\textit{Maîtrise}~= adapte, optimise, forme les autres.

% ==============================================================================
\section{Sujets d'examen types}
\label{sec:sujets-types}
% ==============================================================================

\subsection{Partie~A --- Analyse d'incident (6~points)}

\begin{boxexo}
\textbf{Contexte.}
À 04h15 le 15~mai~2026, le SIEM de l'organisation NovaTech (entreprise
de transfert d'argent mobile, 200~employés, Douala) génère l'alerte
suivante~:

\begin{verbatim}
RULE    : 100205 — Brute force RDP
LEVEL   : 10 (HIGH)
SOURCE  : 196.207.12.88  (ZA, non référencée CTI)
TARGET  : srv-app-01.novatech.local (192.168.20.10)
COUNT   : 12 echecs / 80s  users: Administrator, admin, novatech
MITRE   : T1110.003 — Password Spraying
\end{verbatim}

\textbf{Questions~:}

\begin{enumerate}[label=\textbf{Q\arabic*.}]
  \item (1 pt) Qualifiez l'incident selon la grille de classification.
        Justifiez en 2~phrases.
  \item (1 pt) Listez les 4~premières actions à mener dans les
        15~minutes suivant l'alerte. Indiquez qui les mène (N1 / N2).
  \item (1 pt) Rédigez la commande Wazuh/KQL permettant de vérifier
        si l'IP 196.207.12.88 a réussi une authentification
        RDP (EID~4624, type~10) dans les 2~heures suivant l'alerte.
  \item (2 pts) L'investigation révèle~: connexion RDP réussie à 04h22
        (user: \texttt{administrator}), création d'un nouveau compte
        \texttt{svc\_update} à 04h31 (EID~4720), connexion SMB vers
        srv-bdd-01 à 04h45 (EID~4624). Construisez la timeline annotée
        MITRE ATT\&CK de ces 4~événements (alerte + 3~événements).
  \item (1 pt) Quel artefact forensique, sur le serveur srv-app-01,
        fournirait la preuve d'exécution de commandes par l'attaquant~?
        Comment l'acquérir~?
\end{enumerate}
\end{boxexo}

\subsection{Partie~B --- Investigation forensique (6~points)}

\begin{boxexo}
\textbf{Contexte.}
Vous recevez une image disque \texttt{srv-app-01.dd} et un dump mémoire
\texttt{srv-app-01-ram.lime} acquis à 05h30.

\begin{enumerate}[label=\textbf{Q\arabic*.}]
  \item (1 pt) Décrivez la procédure complète de vérification d'intégrité
        de l'image avant de commencer l'analyse. Quelle commande~?
        Que vérifier~?
  \item (1 pt) Listez 4~commandes Volatility~3 à exécuter en premier
        sur le dump mémoire, dans l'ordre, et justifiez chaque choix.
  \item (2 pts) Le Prefetch révèle l'exécution de \texttt{MSHTA.EXE}
        une seule fois à 04h29. Le registre contient une nouvelle clé
        \texttt{HKCU\textbackslash Run\textbackslash WindowsUpdater}
        pointant vers \texttt{C:\textbackslash ProgramData\textbackslash
        upd.vbs}. Le fichier \texttt{upd.vbs} est absent du disque
        (supprimé). Identifiez les TTPs MITRE correspondants et décrivez
        comment récupérer le fichier supprimé.
  \item (1 pt) Rédigez les 5~éléments obligatoires d'un PV de collecte
        conforme à ISO~27037 pour ce dump mémoire.
  \item (1 pt) La BDD clients (50~000 dossiers) a été accédée entre
        04h45 et 05h15 selon les logs MySQL. Quelles obligations légales
        s'appliquent à NovaTech~? Délai~? Destinataire~?
\end{enumerate}
\end{boxexo}

\subsection{Partie~C --- Supervision et architecture (4~points)}

\begin{boxexo}
\begin{enumerate}[label=\textbf{Q\arabic*.}]
  \item (1 pt) Définissez~: MTTD, MTTR, FPR. Pour chacun~:
        une valeur cible pour un SOC de niveau~3.
  \item (1 pt) NovaTech n'a pas détecté l'attaque RDP à temps
        (MTTD~= 35~jours). Listez 3~sources de logs qui auraient
        dû être collectées et 2~règles SIEM spécifiques au RDP.
  \item (1 pt) Décrivez la différence entre un Playbook et un Runbook.
        Donnez un exemple de chacun applicable à cet incident.
  \item (1 pt) NovaTech souhaite déployer un IDS. Comparez NIDS en mode
        passif et HIDS Wazuh~: pour chaque solution, indiquez ce
        qu'elle aurait détecté ou non dans cet incident.
\end{enumerate}
\end{boxexo}

\subsection{Partie~D --- Cadre légal et RETEX (4~points)}

\begin{boxexo}
\begin{enumerate}[label=\textbf{Q\arabic*.}]
  \item (1 pt) Quelles infractions de la Loi camerounaise~2010/012
        l'attaquant a-t-il commises~? Citez les articles.
  \item (1 pt) NovaTech souhaite porter plainte. Listez les 3~types
        de preuves numériques à constituer, avec leur forme
        (fichier, hash, PV). Quelle est la condition d'admissibilité
        principale~?
  \item (2 pts) Rédigez le plan de remédiation post-incident de NovaTech~:
        6~actions minimum, priorisées (impact, effort, délai,
        responsable, indicateur de vérification).
\end{enumerate}
\end{boxexo}

% ==============================================================================
\section{Corrigé indicatif --- Partie~A}
\label{sec:corrige-partieA}
% ==============================================================================

\begin{boiteRemarque}{Usage pédagogique}
Ce corrigé indicatif est fourni à des fins pédagogiques. D'autres
formulations sont acceptées si elles reflètent une maîtrise correcte
des concepts.
\end{boiteRemarque}

\textbf{Q1 — Qualification.}
Gravité~\gravite{haute} (au minimum)~: l'attaquant cible un serveur
applicatif de production avec du password spraying (T1110.003), technique
plus discrète que le brute force classique. L'escalade vers N2 est
immédiate. Si une authentification réussie est confirmée, escalade vers
\gravite{critique}.

\textbf{Q2 — Premières actions.}
\begin{enumerate}
  \item N1~: Ouvrir ticket INC et documenter l'alerte.
  \item N1~: Vérifier si l'IP est dans la liste blanche
        (partenaire, VPN interne)~; sinon confirmer la qualification.
  \item N2~: Rechercher dans le SIEM si l'IP a réussi une authentification
        dans les 2~h suivantes (EID~4624 type~10).
  \item N2~: Notifier le RSSI et l'Incident Manager~; préparer
        le confinement si succès confirmé.
\end{enumerate}

\textbf{Q3 — Requête KQL.}
\begin{lstlisting}[style=StyleTerminal]
event.code: "4624"
  AND winlog.event_data.LogonType: "10"
  AND winlog.event_data.IpAddress: "196.207.12.88"
  AND @timestamp: [2026-05-15T04:15:00 TO 2026-05-15T06:15:00]
\end{lstlisting}

\textbf{Q4 — Timeline annotée.}

\begin{table}[H]
\centering\small
\rowcolors{2}{TableauPair}{TableauImpair}
\begin{tabular}{p{3cm}p{3.5cm}p{5.5cm}p{2cm}}
\toprule
\tableauHeader{Timestamp} & \tableauHeader{Événement} &
\tableauHeader{Description} & \tableauHeader{MITRE} \\
\midrule
04h15 & EID~\textit{custom} & Password spraying RDP, 12~échecs & T1110.003 \\
04h22 & EID~4624 type~10 & Auth. RDP réussie (administrator) & T1078 \\
04h31 & EID~4720 & Création compte svc\_update & T1136.001 \\
04h45 & EID~4624 type~3 & Connexion SMB srv-bdd-01 & T1021.002 \\
\bottomrule
\end{tabular}
\end{table}

\textbf{Q5 — Artefact forensique.}
Le \textbf{Prefetch} (\texttt{C:\textbackslash Windows\textbackslash Prefetch\textbackslash}) fournirait la preuve d'exécution~:
nom de l'exécutable, date/heure de dernière exécution, nombre d'exécutions,
DLLs chargées. Acquisition~: copie depuis une image montée en lecture seule~;
ou \texttt{PECmd.exe -d C:\textbackslash Windows\textbackslash Prefetch --csv output/}.

% ==============================================================================
\section{Barème des laboratoires}
\label{sec:bareme-labs}
% ==============================================================================

\begin{table}[H]
\centering\small
\caption{Barème des rapports de laboratoires}
\label{tab:bareme-labs}
\rowcolors{2}{TableauPair}{TableauImpair}
\begin{tabular}{p{1cm}p{4cm}p{6.5cm}p{2cm}}
\toprule
\tableauHeader{Lab} & \tableauHeader{Titre} &
\tableauHeader{Critères d'évaluation} & \tableauHeader{Coefficient} \\
\midrule
1 & SIEM et détection &
Configuration correcte du pipeline (4 pts)~;
règles opérationnelles (3 pts)~;
qualité du tableau de bord (2 pts)~;
rapport d'alerte structuré (1 pt) &
$\times$~1 \\
2 & Forensique disque &
PV de collecte conforme (3 pts)~;
timeline correcte avec MITRE (4 pts)~;
artefacts identifiés (2 pts)~;
qualité du rapport (1 pt) &
$\times$~1.5 \\
3 & Simulation Tabletop ransomware &
Qualification de l'incident (2 pts)~;
playbook appliqué (3 pts)~;
plan de remédiation (3 pts)~;
rapport exécutif (2 pts) &
$\times$~1 \\
\bottomrule
\end{tabular}
\end{table}

% ==============================================================================
\section{Conseils méthodologiques pour l'examen}
\label{sec:conseils-exam}
% ==============================================================================

\begin{boxinfo}
\textbf{Pour l'examen final~:}
\begin{itemize}
  \item \textbf{Lire le sujet en entier avant d'écrire.} Les questions
        s'enchaînent~: la réponse à Q3 peut dépendre de Q2.
  \item \textbf{Être précis sur les noms d'outils et les EID.}
        \og{}Un log Windows\fg{} n'est pas une réponse~;
        \og{}Event ID~4624 dans le journal Security\fg{} en est une.
  \item \textbf{Citer les références normatives.} Pour toute question
        sur les preuves ou les obligations légales, citer la norme
        (ISO~27037, RFC~3227, Loi~2010/012 + article).
  \item \textbf{Structurer les réponses techniques.} Pour une timeline~:
        tableau avec colonnes Timestamp / Source / Événement / MITRE.
        Pour un playbook~: étapes numérotées avec responsable.
  \item \textbf{Faire preuve de bon sens opérationnel.} Avant de
        répondre techniquement, se demander~: \og{}Dans la vraie vie,
        qu'est-ce qui est prioritaire ici~?\fg{}
\end{itemize}
\end{boxinfo}

% ==============================================================================
\section*{Récapitulatif des compétences du cours}
% ==============================================================================

\begin{boxresume}
\textbf{À l'issue de SIS589, vous êtes capable de~:}
\begin{itemize}
  \item Opérer un SOC minimal~: déployer Wazuh + ELK, rédiger des règles
        de corrélation SIEM, calculer et interpréter MTTD / MTTR / FPR.
  \item Lire, corréler et interpréter des journaux hétérogènes pour
        reconstruire la chronologie d'une attaque réelle.
  \item Acquérir des preuves numériques (disque, RAM, réseau) dans le
        respect des principes forensiques (intégrité, chaîne de custody,
        ISO~27037 / RFC~3227).
  \item Analyser les artefacts Windows (MFT, Prefetch, registre, Sysmon)
        et Linux (inodes, bash\_history, persistances) en utilisant
        Volatility~3, The Sleuth Kit, Autopsy.
  \item Piloter une réponse à incident selon ISO~27035~: qualification,
        confinement, éradication, RETEX.
  \item Appliquer le cadre juridique camerounais (Loi~2010/012) et
        produire des preuves légalement recevables.
  \item Intégrer la Cyber Threat Intelligence (MISP, STIX, ATT\&CK)
        dans les processus de détection et de réponse.
  \item Conduire un RETEX structuré et un plan de remédiation priorisé.
\end{itemize}
\end{boxresume}
